\chapter{Stoffe und Stoffeigenschaften}

\section{Schematischer Aufbau der Atome}

Atome bestehen aus dem positiv geladenem
Kern und der negativen Elektronenhülle.

\subsection{Kern}

Der Kern setzt sich aus einzelnen Teilchen zu-
sammen: den positiv geladenen Protonen und den
ungeladenen Neutronen. Die Masse von Proton
und Neutron ist nahezu gleich.

Der Kernradius ist ausserordentlich klein
($r_K = 10^{-15} m$). Er hängt von der Zahl der Protonen 
und Neutronen ab, die man Nukleonenzahl oder 
Massenzahl A nennt. Die Nukleonenzahl A ist die Summe aus Protonenzahl Z (DIN
32640) und Neutronenzahl N \cite{boge_physik_2008}.

\begin{minipage}{.5\textwidth}
    \begin{equbox}
        \begin{equation}
            A = Z + N
        \end{equation}
    \end{equbox}
\end{minipage}%
\begin{minipage}{.5\textwidth}
    % Image on the right
    \centering
    \captionsetup{type=figure}
    \includegraphics[clip, trim=0cm 0cm 0cm 0cm, width=.7\textwidth]{aufbau_atomkern.png}
    \captionof{figure}{Aufbau des Atomkerns \cite{boge_physik_2008}}
    \label{fig:aufbauAtomkern}
\end{minipage}

\subsection{Atomhülle}

Die Atomhülle ist etwa 100000 mal ausgedehnter 
als der Kern. Die Atomradien liegen bei
etwa $10^{-10}m$. Die Hülle wird durch Elektronen
gebildet, die etwa 2000 mal leichter als Protonen 
oder Neutronen sind \cite{boge_physik_2008}.
Zahl der Elektronen in der Hülle =
Zahl der Protonen im Kern.

\subsection{Atommodell}

Aus den obigen Daten erhält man folgendes
Bild: Die Masse eines Atoms ist zu über 99,9\%
im sehr kleinen Kern konzentriert. Die leichten
Elektronen umkreisen den winzigen schweren
Kern in relativ grossen Abständen. Nach diesen
Vorstellungen besteht das Atom aus den Elektronenbahnen 
und dem Kern. Dazwischen ist
leerer Raum.

Alle chemischen und elektrischen Prozesse der
Materie finden in der Atomhülle statt. Das Gleiche 
gilt für mechanische Vorgänge, die auf den
Kräften zwischen den Elektronenhüllen beruhen. 
Der Kern ist durch die Hülle abgeschirmt.
Seine positive Ladung hält das Atom zusammen \cite{boge_physik_2008}.

\begin{minipage}{.6\textwidth}
    % Image on the right
    \centering
    \captionsetup{type=figure}
    \includegraphics[clip, trim=0cm 0cm 0cm 0cm, width=.7\textwidth]{atommodell.png}
    \captionof{figure}{Das gängigste Atommodell \cite{boge_physik_2008}}
    \label{fig:atommodell}
\end{minipage}%
\begin{minipage}{.4\textwidth}
    \begin{infobox}
        Chemische und elektrische Vorgänge laufen 
        in der Atomhülle ab.
    \end{infobox}
\end{minipage}

\subsection{Periodensystem}

\begin{infobox}
    Jedes Element wird durch die Ordnungszahl
    Z (und die Massenzahl A) charakterisiert.
    \begin{itemize}
        \item Ordnungszahl $Z$ = Zahl der Elektronen in der Hülle, beschreibt
            die Stellung des Atoms im Periodensystem. 
    \end{itemize}
\end{infobox}

Das leichteste Atom ist Wasserstoff H mit $Z = 1$.
Es besteht aus einem Proton im Kern und einem Elektron in der Atomhülle.
Das nächste Atom ist Helium He mit $Z = 2$. Im
Kern befinden sich zwei Protonen und in der
Hülle zwei Elektronen. Abbildung \ref{fig:periodensystem}
zeigt das Periodensystem (mit Ordnungszahl) für die meisten bekannten Stoffe.

\begin{figure}[H]    
    \centering          %     left botm right top
    \includegraphics[width=.9\textwidth]{periodensystem.png}
    \caption{Periodensystem der Elemente \cite{noauthor_periodensystem_nodate}}
    \label{fig:periodensystem}
\end{figure}

Das Wasserstoffatom hat in seiner Atomhülle
nur ein Elektron. Alle weiteren Atome besitzen
mehrere Elektronen, die sich in so genannte
Schalen gruppieren. Jede Schale kann nur eine
bestimmte Anzahl von Elektronen aufnehmen
(Pauli-Prinzip). Man kann alle Atome in ein
Periodensystem einordnen, das den Schalenaufbau widerspiegelt.

Jede Periode stellt eine gefüllte Elektronenschale dar. Dementsprechend hat ein Atom,
das am Anfang einer Periode steht, ein Elektron
ausserhalb einer abgeschlossenen Schale (Alkalien).
Ein Atom am Ende einer Periode besitzt
eine abgeschlossene Schale und ist daher
chemisch besonders stabil (Edelgas) \cite{boge_physik_2008}.

\begin{itemize}
    \item 1. Periode: H und He
    \item 2. Periode: Li, Be, B, C, N, O, F, Ne
    \item 3. Periode: Na, Mg, Al, Si, P, S. Cl, Ar
    \item usw.
\end{itemize}

\subsection{Isotop, Molekül, Ion}

Ein Isotop ist eine Variante eines chemischen Elements, die sich 
durch die Anzahl der Neutronen im Atomkern unterscheidet, 
während die Protonenzahl (und damit die chemischen Eigenschaften) gleich bleibt.
An Wasserstoff als Beispiel, Abbildung \ref{fig:isotopedeswasserstoffs}.

\begin{figure}[H]    
    \centering          %     left botm right top
    \includegraphics[width=.6\textwidth]{isotop_schematisch_wasserstoff.png}
    \caption{Isotope des Wasserstoffs \cite{noauthor_isotope_nodate}}
    \label{fig:isotopedeswasserstoffs}
\end{figure}

Ein Molekül besteht aus zwei oder mehr Atomen, die durch chemische Bindungen 
miteinander verbunden sind und als Einheit agieren.

\begin{figure}[H]    
    \centering          %     left botm right top
    \includegraphics[width=.5\textwidth]{molekül_schematisch_beispiel.png}
    \caption{Beispiel eines Moleküls \cite{noauthor_pfas_2023}}
    \label{fig:molekülschematisch}
\end{figure}

Ein Ion ist ein Atom oder Molekül, das durch Gewinn oder Verlust von 
Elektronen elektrisch geladen ist (positiv als Kation, negativ als Anion).

\begin{figure}[H]    
    \centering          %     left botm right top
    \includegraphics[width=.5\textwidth]{ion_schematisch.jpg}
    \caption{Beispiel eines Moleküls \cite{noauthor_was_2022}}
    \label{fig:ionschematisch}
\end{figure}

\section{Aggregatszustände}

Stoffe können in unterschiedlichen Aggregatzuständen vorkommen. 
Die physikalischen Eigenschaften ändern sich, das Element bleibt dabei dasselbe.

In welchem Aggregatzustand sich ein Stoff befindet, hängt von der 
Temperatur und dem Druck ab. Dabei ist die Temperatur eine statistische 
Grösse, sie beschreibt die mittlere kinetische Energie aller Teilchen. 
Die physikalischen Auswirkungen von Temperaturänderungen werden durch die 
Wärmelehre beschrieben \cite{noauthor_1_nodate}.

\subsection{Phasendiagramm}

Um den Zusammenhang zwischen Aggregatzustand, Druck und 
Temperatur darzustellen, werden Phasendiagramme verwendet.
Abbildung \ref{fig:phasendiagrammwasser} zeigt ein Phasendiagramm (mit Anomalie)
am Beispiel Wasser.

\begin{infobox}
    Der \textbf{Tripelpunkt} ein Zustand eines aus einer einzigen 
    Stoffkomponente bestehenden Systems, in dem Temperatur und Druck dreier 
    Phasen im thermodynamischen Gleichgewicht stehen.
\end{infobox}

\begin{infobox}
    Der \textbf{kritische Punkt} ist ein thermodynamischer Zustand 
    eines Stoffes, der sich durch Angleichen der Dichte von flüssiger und 
    gasförmiger Phase kennzeichnet. Die Unterschiede zwischen beiden 
    Aggregatzuständen hören an diesem Punkt auf zu existieren.
\end{infobox}

\begin{figure}[H]
    \centering          %     left botm right top
    \includegraphics[width=.7\textwidth]{phasendiagramm_wasser.png}
    \caption{Phasendiagramm Wasser \cite{noauthor_1_nodate}}
    \label{fig:phasendiagrammwasser}
\end{figure}

\subsection{Übergänge}

Konkret bedeutet das, dass sich die Aggregatzustände in Abhängigkeit
von Aussenbedingungen ändern können. Eine Wärmezufuhr bringt feste Stoffe 
zum Schmelzen, und Flüssigkeiten
können durch Druckerniedrigung oder Temperaturerhöhung verdampfen.
Beim Übergang von einem zum anderen Aggregatzustand erfolgt der
Wechsel mit steigender Temperatur nicht allmählich und kontinuierlich,
sondern plötzlich und sprunghaft.
Es gibt vier verschiedene Übergänge, wobei die Temperatur, 
bei der ein Aggregatzustand in einen anderen übergeht, für jeden Stoff
anders und für jeden auch eine charakteristische Grösse ist \cite{bannwarth_basiswissen_2007-1}.

\begin{infobox}
    \begin{minipage}{.3\textwidth}
        \textbf{Fixpunkt} \\
        Siedepunkt \\
        Kondensationspunkt \\
        Schmelzpunkt \\
        Erstarrungspunkt
    \end{minipage}
    \hspace{.1\textwidth}
    \begin{minipage}{.3\textwidth}
        \textbf{Übergang} \\
        flüssig    ->    gasförmig \\
        gasförmig  ->    flüssig \\
        fest       ->    flüssig \\
        flüssig    ->    fest  
    \end{minipage}
\end{infobox}